% !TEX program = lualatex
\documentclass[aspectratio=43,professionalfonts,handout]{beamer}
\usepackage{beamer_template}

\newcommand{\LectureNum}{7}
\newcommand{\LectureTitle}{線形システムの伝達関数とデルタ関数}
\newcommand{\TextPages}{p.\,68-71}
\newcommand{\CourseName}{応用数学}
\renewcommand{\TermName}{前期}

\begin{document}

\begin{frame}{\CourseName\quad 第\LectureNum 回「\LectureTitle」}
  {\small \TermName\quad \TeacherName\quad ／\quad 教科書 \TextPages}

  \textbf{目標}：線形システムの伝達関数とデルタ関数の概念を理解できる

\end{frame}

\begin{frame}{線形システム}
  \begin{block}{}
  \begin{itemize}
    \item 入力 $x(t) \to $ システム $ \to $ 出力 $y(t)$
    \item 伝達関数 $H(s) = Y(s)/X(s)$
    \item $y(t) = h(t)*x(t)$
  \end{itemize}
  \end{block}
\end{frame}

\begin{frame}{伝達関数}
  \begin{block}{}
  \begin{itemize}
    \item 微分方程式 $ \to $ ラプラス変換 $ \to H(s)$
    \item $H(s) = 1/(s^{2}+as+b)$
    \item $h(t) = \mathcal{L}^{-1}[H(s)]$（インパルス応答）
  \end{itemize}
  \end{block}
\end{frame}

\begin{frame}{デルタ関数}
  \begin{block}{}
  \begin{itemize}
    \item $\varphi_\varepsilon(t)$ の極限としてのデルタ関数
    \item $\mathcal{L}[\delta(t)] = 1$
    \item $\int\delta(t)dt = 1$
    \item $f*\delta = f$（たたみ込みとの関係）
  \end{itemize}
  \end{block}
\end{frame}

\begin{frame}{演習（練習問題2）}
  \begin{exampleblock}{自分で解いてみよう}
  \begin{itemize}
    \item p.72 問1〜5より選題
  \end{itemize}
  \end{exampleblock}
\end{frame}

\begin{frame}{ラプラス変換まとめ}
  \begin{itemize}
    \item 第1〜7週の総括
    \item 前期中間試験に向けて
    \item 次回：フーリエ級数
  \end{itemize}
\end{frame}

\end{document}
